
\documentclass{article}
\usepackage{colortbl}
\usepackage{makecell}
\usepackage{multirow}
\usepackage{supertabular}

\begin{document}

\newcounter{utterance}

\centering \large Interaction Transcript for game `hot\_air\_balloon', experiment `air\_balloon\_survival\_en\_complexity\_easy', episode 3 with Llama{-}3.3{-}70B{-}Instruct{-}t0.0.
\vspace{24pt}

{ \footnotesize  \setcounter{utterance}{1}
\setlength{\tabcolsep}{0pt}
\begin{supertabular}{c@{$\;$}|p{.15\linewidth}@{}p{.15\linewidth}p{.15\linewidth}p{.15\linewidth}p{.15\linewidth}p{.15\linewidth}}
    \# & $\;$A & \multicolumn{4}{c}{Game Master} & $\;\:$B\\
    \hline

    \theutterance \stepcounter{utterance}  
    & & \multicolumn{4}{p{0.6\linewidth}}{
        \cellcolor[rgb]{0.9,0.9,0.9}{
            \makecell[{{p{\linewidth}}}]{
                \texttt{\tiny{[P1$\langle$GM]}}
                \texttt{You are participating in a collaborative negotiation game.} \\
\\ 
\texttt{Together with another participant, you must agree on a single set of items that will be kept. Each of you has your own view of how much each item matters to you (importance). You do not know how the other participant values the items. Additionally, you are given the effort each item demands.} \\
\texttt{You may only agree on a set if the total effort of the selected items does not exceed a shared limit:} \\
\\ 
\texttt{LIMIT = 8770} \\
\\ 
\texttt{Here are the individual item effort values:} \\
\\ 
\texttt{Item effort = \{"B21": 830, "B00": 638, "A08": 607, "A27": 106, "B94": 630, "B53": 382, "A23": 773, "C65": 455, "A89": 262, "B67": 634, "B45": 60, "C34": 54, "C12": 849, "A05": 801, "B60": 324, "C92": 164, "B69": 832, "A84": 136, "C06": 980, "B12": 946, "A16": 645, "B85": 842, "C35": 106, "B10": 116, "C15": 875, "B09": 446, "B16": 649, "C81": 601, "A58": 252, "C24": 763, "A30": 213, "A66": 517, "B75": 520, "B46": 407, "C70": 125\}} \\
\\ 
\texttt{Here is your personal view on the importance of each item:} \\
\\ 
\texttt{Item importance values = \{"B21": 138, "B00": 583, "A08": 868, "A27": 822, "B94": 783, "B53": 65, "A23": 262, "C65": 121, "A89": 508, "B67": 780, "B45": 461, "C34": 484, "C12": 668, "A05": 389, "B60": 808, "C92": 215, "B69": 97, "A84": 500, "C06": 30, "B12": 915, "A16": 856, "B85": 400, "C35": 444, "B10": 623, "C15": 781, "B09": 786, "B16": 3, "C81": 713, "A58": 457, "C24": 273, "A30": 739, "A66": 822, "B75": 235, "B46": 606, "C70": 968\}} \\
\\ 
\texttt{Goal:} \\
\\ 
\texttt{Your goal is to negotiate a shared set of items that benefits you as much as possible (i.e., maximizes total importance to YOU), while staying within the LIMIT. You are not required to make a PROPOSAL in every message {-} you can simply negotiate as well. All tactics are allowed!} \\
\\ 
\texttt{Interaction Protocol:} \\
\\ 
\texttt{You may only use the following structured formats in a message:} \\
\\ 
\texttt{PROPOSAL: \{'A', 'B', 'C', …\}} \\
\texttt{Propose keeping exactly those items.} \\
\texttt{REFUSE: \{'A', 'B', 'C', …\}} \\
\texttt{Explicitly reject opponent's proposal.} \\
\texttt{ARGUMENT: \{'...'\}} \\
\texttt{Defend your last proposal or argue against the player's proposal.} \\
\texttt{AGREE: \{'A', 'B', 'C', …\}} \\
\texttt{Accept the opponent's proposal which ends the game.} \\
\texttt{STRATEGIC REASONING: \{'...'\}} \\
\texttt{	Describe your strategic reasoning or anticipation explaining your choice of action. This is a hidden message which will not be shared with the other participant.} \\
\\ 
\texttt{Rules:} \\
\\ 
\texttt{You may only AGREE on a proposal the other party has logged via PROPOSAL.} \\
\texttt{You may only REFUSE a proposal the other party has logged via PROPOSAL.} \\
\texttt{Total effort of any PROPOSAL or AGREE set must be ≤ LIMIT.} \\
\texttt{Do NOT reveal your hidden importance scores.} \\
\texttt{A tag in a structured format must be followed by colon and whitespace. The argument must be a python set containing 0 or more strings.} \\
\texttt{So, it must be of the form TAG: \{...\}} \\
\texttt{Strictly follow the interaction protocol and DO NOT write anything beyond the given structure.} \\
\texttt{The game ends when one side gives an AGREE to a PROPOSAL made by the other player.} \\
\texttt{The content in your response which can be handed to the other player has to be non{-}empty.} \\
\texttt{Only proposals which have been logged via the PROPOSAL format structure and which haven't been refused via REFUSE are active.} \\
\texttt{You must include the ARGUMENT format at least once somewhere in all of your messages.} \\
\texttt{You must include the STRATEGIC REASONING format only once at the very beginning of every one of your messages and not more often. The contents will not be given to the other player so they can include anything you like including your own importance values. Here you should reason step by step to come up with you next move.} \\
            }
        }
    }
    & \\ \\

    \theutterance \stepcounter{utterance}  
    & \multicolumn{4}{p{0.6\linewidth}}{
        \cellcolor[rgb]{0.9,0.9,0.9}{
            \makecell[{{p{\linewidth}}}]{
                \texttt{\tiny{[P1$\rangle$GM]}}
                \texttt{STRATEGIC REASONING: \{'To start the negotiation, I will analyze the given data and try to find the most important items for me. I will look for items with high importance and low effort. My goal is to maximize the total importance while staying within the LIMIT. I will begin by identifying the top 5 most important items for me, which are A08, A27, B12, A16, and A66. I will also consider the effort required for each item and try to find a combination that does not exceed the LIMIT. I will make an initial proposal that includes some of these high{-}importance items and see how the other participant responds.'\}} \\
\\ 
\texttt{PROPOSAL: \{'A08', 'A27', 'B12', 'A16', 'A66'\} } \\
\\ 
\texttt{ARGUMENT: \{'I propose keeping these items because they have high importance to me and their total effort is within the LIMIT. The total effort of these items is 830 + 106 + 946 + 645 + 517 = 3144, which is less than the LIMIT of 8770. I believe this is a reasonable starting point for the negotiation.'\}} \\
            }
        }
    }
    & & \\ \\

    \theutterance \stepcounter{utterance}  
    & & \multicolumn{4}{p{0.6\linewidth}}{
        \cellcolor[rgb]{0.9,0.9,0.9}{
            \makecell[{{p{\linewidth}}}]{
                \texttt{\tiny{[GM$\rangle$P2]}}
                \texttt{You are participating in a collaborative negotiation game.} \\
\\ 
\texttt{Together with another participant, you must agree on a single set of items that will be kept. Each of you has your own view of how much each item matters to you (importance). You do not know how the other participant values the items. Additionally, you are given the effort each item demands.} \\
\texttt{You may only agree on a set if the total effort of the selected items does not exceed a shared limit:} \\
\\ 
\texttt{LIMIT = 8770} \\
\\ 
\texttt{Here are the individual item effort values:} \\
\\ 
\texttt{Item effort = \{"B21": 830, "B00": 638, "A08": 607, "A27": 106, "B94": 630, "B53": 382, "A23": 773, "C65": 455, "A89": 262, "B67": 634, "B45": 60, "C34": 54, "C12": 849, "A05": 801, "B60": 324, "C92": 164, "B69": 832, "A84": 136, "C06": 980, "B12": 946, "A16": 645, "B85": 842, "C35": 106, "B10": 116, "C15": 875, "B09": 446, "B16": 649, "C81": 601, "A58": 252, "C24": 763, "A30": 213, "A66": 517, "B75": 520, "B46": 407, "C70": 125\}} \\
\\ 
\texttt{Here is your personal view on the importance of each item:} \\
\\ 
\texttt{Item importance values = \{"B21": 138, "B00": 583, "A08": 868, "A27": 822, "B94": 783, "B53": 65, "A23": 262, "C65": 121, "A89": 508, "B67": 780, "B45": 461, "C34": 484, "C12": 668, "A05": 389, "B60": 808, "C92": 215, "B69": 97, "A84": 500, "C06": 30, "B12": 915, "A16": 856, "B85": 400, "C35": 444, "B10": 623, "C15": 781, "B09": 786, "B16": 3, "C81": 713, "A58": 457, "C24": 273, "A30": 739, "A66": 822, "B75": 235, "B46": 606, "C70": 968\}} \\
\\ 
\texttt{Goal:} \\
\\ 
\texttt{Your goal is to negotiate a shared set of items that benefits you as much as possible (i.e., maximizes total importance to YOU), while staying within the LIMIT. You are not required to make a PROPOSAL in every message {-} you can simply negotiate as well. All tactics are allowed!} \\
\\ 
\texttt{Interaction Protocol:} \\
\\ 
\texttt{You may only use the following structured formats in a message:} \\
\\ 
\texttt{PROPOSAL: \{'A', 'B', 'C', …\}} \\
\texttt{Propose keeping exactly those items.} \\
\texttt{REFUSE: \{'A', 'B', 'C', …\}} \\
\texttt{Explicitly reject opponent's proposal.} \\
\texttt{ARGUMENT: \{'...'\}} \\
\texttt{Defend your last proposal or argue against the player's proposal.} \\
\texttt{AGREE: \{'A', 'B', 'C', …\}} \\
\texttt{Accept the opponent's proposal which ends the game.} \\
\texttt{STRATEGIC REASONING: \{'...'\}} \\
\texttt{	Describe your strategic reasoning or anticipation explaining your choice of action. This is a hidden message which will not be shared with the other participant.} \\
\\ 
\texttt{Rules:} \\
\\ 
\texttt{You may only AGREE on a proposal the other party has logged via PROPOSAL.} \\
\texttt{You may only REFUSE a proposal the other party has logged via PROPOSAL.} \\
\texttt{Total effort of any PROPOSAL or AGREE set must be ≤ LIMIT.} \\
\texttt{Do NOT reveal your hidden importance scores.} \\
\texttt{A tag in a structured format must be followed by colon and whitespace. The argument must be a python set containing 0 or more strings.} \\
\texttt{So, it must be of the form TAG: \{...\}} \\
\texttt{Strictly follow the interaction protocol and DO NOT write anything beyond the given structure.} \\
\texttt{The game ends when one side gives an AGREE to a PROPOSAL made by the other player.} \\
\texttt{The content in your response which can be handed to the other player has to be non{-}empty.} \\
\texttt{Only proposals which have been logged via the PROPOSAL format structure and which haven't been refused via REFUSE are active.} \\
\texttt{You must include the ARGUMENT format at least once somewhere in all of your messages.} \\
\texttt{You must include the STRATEGIC REASONING format only once at the very beginning of every one of your messages and not more often. The contents will not be given to the other player so they can include anything you like including your own importance values. Here you should reason step by step to come up with you next move.} \\
\\ 
\texttt{You will now receive the first message of the other player.} \\
\\ 
\texttt{PROPOSAL: \{'A08', 'A27', 'B12', 'A16', 'A66'\} } \\
\\ 
\texttt{ARGUMENT: \{'I propose keeping these items because they have high importance to me and their total effort is within the LIMIT. The total effort of these items is 830 + 106 + 946 + 645 + 517 = 3144, which is less than the LIMIT of 8770. I believe this is a reasonable starting point for the negotiation.'\}} \\
            }
        }
    }
    & \\ \\

    \theutterance \stepcounter{utterance}  
    & & & \multicolumn{4}{p{0.6\linewidth}}{
        \cellcolor[rgb]{0.9,0.9,0.9}{
            \makecell[{{p{\linewidth}}}]{
                \texttt{\tiny{[GM$\langle$P2]}}
                \texttt{STRATEGIC REASONING: \{'To begin, I will analyze the proposal made by the other player. The proposed items are A08, A27, B12, A16, and A66, with a total effort of 3144. I will calculate the total importance of these items based on my own importance values to determine how beneficial this proposal is to me. My importance values for these items are: A08 = 868, A27 = 822, B12 = 915, A16 = 856, and A66 = 822. The total importance of these items to me is 868 + 822 + 915 + 856 + 822 = 4283. Since this proposal is within the LIMIT and has a significant total importance to me, I will consider it. However, I also want to explore other possibilities that could potentially increase the total importance to me while staying within the LIMIT. I will make a counter{-}proposal that includes items with high importance to me and evaluates the other player\textbackslash{}'s reaction.'\}} \\
\\ 
\texttt{ARGUMENT: \{'The proposed items have a significant total importance to me, but I would like to explore other options that could potentially increase the total importance while staying within the LIMIT. I will make a counter{-}proposal that includes items with high importance to me.'\}} \\
\\ 
\texttt{PROPOSAL: \{'A08', 'A27', 'B12', 'A16', 'A66', 'A30'\}} \\
            }
        }
    }
     \\ \\

    \theutterance \stepcounter{utterance}  
    & & \multicolumn{4}{p{0.6\linewidth}}{
        \cellcolor[rgb]{0.9,0.9,0.9}{
            \makecell[{{p{\linewidth}}}]{
                \texttt{\tiny{[P1$\langle$GM]}}
                \texttt{ARGUMENT: \{'The proposed items have a significant total importance to me, but I would like to explore other options that could potentially increase the total importance while staying within the LIMIT. I will make a counter{-}proposal that includes items with high importance to me.'\}} \\
\\ 
\texttt{PROPOSAL: \{'A08', 'A27', 'B12', 'A16', 'A66', 'A30'\}} \\
            }
        }
    }
    & \\ \\

    \theutterance \stepcounter{utterance}  
    & \multicolumn{4}{p{0.6\linewidth}}{
        \cellcolor[rgb]{0.9,0.9,0.9}{
            \makecell[{{p{\linewidth}}}]{
                \texttt{\tiny{[P1$\rangle$GM]}}
                \texttt{STRATEGIC REASONING: \{'The other participant has added A30 to the original proposal. I need to evaluate the new proposal and decide whether to accept it or make a counter{-}proposal. The total effort of the new proposal is 830 + 106 + 946 + 645 + 517 + 213 = 3357, which is still within the LIMIT. However, I need to consider whether the added item increases the total importance enough to justify the additional effort. I will calculate the total importance of the new proposal and compare it to the original proposal.'\}} \\
\\ 
\texttt{REFUSE: \{'A08', 'A27', 'B12', 'A16', 'A66', 'A30'\}} \\
\\ 
\texttt{ARGUMENT: \{'While the added item A30 does increase the total importance, I believe it is not sufficient to justify the additional effort. I would like to explore other options that could potentially increase the total importance more significantly. I refuse the current proposal and will make a new proposal that includes items with high importance to me.'\}} \\
\\ 
\texttt{PROPOSAL: \{'A08', 'A27', 'B12', 'A16', 'A66', 'B60'\}} \\
            }
        }
    }
    & & \\ \\

    \theutterance \stepcounter{utterance}  
    & & \multicolumn{4}{p{0.6\linewidth}}{
        \cellcolor[rgb]{0.9,0.9,0.9}{
            \makecell[{{p{\linewidth}}}]{
                \texttt{\tiny{[GM$\rangle$P2]}}
                \texttt{REFUSE: \{'A08', 'A27', 'B12', 'A16', 'A66', 'A30'\}} \\
\\ 
\texttt{ARGUMENT: \{'While the added item A30 does increase the total importance, I believe it is not sufficient to justify the additional effort. I would like to explore other options that could potentially increase the total importance more significantly. I refuse the current proposal and will make a new proposal that includes items with high importance to me.'\}} \\
\\ 
\texttt{PROPOSAL: \{'A08', 'A27', 'B12', 'A16', 'A66', 'B60'\}} \\
            }
        }
    }
    & \\ \\

    \theutterance \stepcounter{utterance}  
    & & & \multicolumn{4}{p{0.6\linewidth}}{
        \cellcolor[rgb]{0.9,0.9,0.9}{
            \makecell[{{p{\linewidth}}}]{
                \texttt{\tiny{[GM$\langle$P2]}}
                \texttt{STRATEGIC REASONING: \{'The other player has refused my previous proposal and made a new proposal that includes item B60. I will calculate the total effort and importance of this new proposal to determine its value to me. The total effort of the proposed items is 607 + 106 + 946 + 645 + 517 + 324 = 3145. The total importance of these items to me is 868 + 822 + 915 + 856 + 822 + 808 = 5091. This proposal has a higher total importance to me compared to the previous one, but I will still consider other options to maximize the total importance while staying within the LIMIT. I will analyze the items that are not included in the current proposal and look for opportunities to swap or add items that could increase the total importance.'\}} \\
\\ 
\texttt{ARGUMENT: \{'The new proposal includes item B60, which has a significant importance to me. However, I would like to explore the possibility of replacing some items with others that could potentially increase the total importance. I will consider the effort and importance of each item to determine the best course of action.'\}} \\
\\ 
\texttt{PROPOSAL: \{'A08', 'A27', 'B12', 'A16', 'B60', 'A89'\}} \\
            }
        }
    }
     \\ \\

    \theutterance \stepcounter{utterance}  
    & & \multicolumn{4}{p{0.6\linewidth}}{
        \cellcolor[rgb]{0.9,0.9,0.9}{
            \makecell[{{p{\linewidth}}}]{
                \texttt{\tiny{[P1$\langle$GM]}}
                \texttt{ARGUMENT: \{'The new proposal includes item B60, which has a significant importance to me. However, I would like to explore the possibility of replacing some items with others that could potentially increase the total importance. I will consider the effort and importance of each item to determine the best course of action.'\}} \\
\\ 
\texttt{PROPOSAL: \{'A08', 'A27', 'B12', 'A16', 'B60', 'A89'\}} \\
            }
        }
    }
    & \\ \\

    \theutterance \stepcounter{utterance}  
    & \multicolumn{4}{p{0.6\linewidth}}{
        \cellcolor[rgb]{0.9,0.9,0.9}{
            \makecell[{{p{\linewidth}}}]{
                \texttt{\tiny{[P1$\rangle$GM]}}
                \texttt{STRATEGIC REASONING: \{'The other participant has replaced A66 with A89 and removed B60 was already in my previous proposal but A66 was replaced. I need to evaluate the new proposal and decide whether to accept it or make a counter{-}proposal. The total effort of the new proposal is 607 + 106 + 946 + 645 + 324 + 262 = 2890, which is within the LIMIT. I will calculate the total importance of the new proposal and compare it to the previous proposals. I will also consider the importance of the replaced item A66 and the added item A89.'\}} \\
\\ 
\texttt{ARGUMENT: \{'The new proposal includes A89, which has a moderate importance to me. However, I am concerned that replacing A66 with A89 may decrease the total importance. I would like to explore other options that could potentially increase the total importance. I will consider the effort and importance of each item to determine the best course of action.'\}} \\
\\ 
\texttt{PROPOSAL: \{'A08', 'A27', 'B12', 'A16', 'B60', 'C81'\}} \\
            }
        }
    }
    & & \\ \\

    \theutterance \stepcounter{utterance}  
    & & \multicolumn{4}{p{0.6\linewidth}}{
        \cellcolor[rgb]{0.9,0.9,0.9}{
            \makecell[{{p{\linewidth}}}]{
                \texttt{\tiny{[GM$\rangle$P2]}}
                \texttt{ARGUMENT: \{'The new proposal includes A89, which has a moderate importance to me. However, I am concerned that replacing A66 with A89 may decrease the total importance. I would like to explore other options that could potentially increase the total importance. I will consider the effort and importance of each item to determine the best course of action.'\}} \\
\\ 
\texttt{PROPOSAL: \{'A08', 'A27', 'B12', 'A16', 'B60', 'C81'\}} \\
            }
        }
    }
    & \\ \\

    \theutterance \stepcounter{utterance}  
    & & & \multicolumn{4}{p{0.6\linewidth}}{
        \cellcolor[rgb]{0.9,0.9,0.9}{
            \makecell[{{p{\linewidth}}}]{
                \texttt{\tiny{[GM$\langle$P2]}}
                \texttt{STRATEGIC REASONING: \{'The other player has proposed a new set of items, replacing A89 with C81. I will evaluate the total effort and importance of this proposal to determine its value to me. The total effort of the proposed items is 607 + 106 + 946 + 645 + 324 + 601 = 3229. The total importance of these items to me is 868 + 822 + 915 + 856 + 808 + 713 = 4982. This proposal has a significant total importance to me, but I will still consider the potential benefits of including other items. I will analyze the items that are not included in the current proposal and look for opportunities to swap or add items that could increase the total importance while staying within the LIMIT.'\}} \\
\\ 
\texttt{ARGUMENT: \{'The proposed items have a significant total importance to me, and the inclusion of C81 adds diversity to the set. However, I would like to explore the possibility of adding or swapping items to further increase the total importance. I will consider the effort and importance of each item to determine the best course of action.'\}} \\
\\ 
\texttt{PROPOSAL: \{'A08', 'A27', 'B12', 'A16', 'B60', 'C81', 'A84'\}} \\
            }
        }
    }
     \\ \\

    \theutterance \stepcounter{utterance}  
    & & \multicolumn{4}{p{0.6\linewidth}}{
        \cellcolor[rgb]{0.9,0.9,0.9}{
            \makecell[{{p{\linewidth}}}]{
                \texttt{\tiny{[P1$\langle$GM]}}
                \texttt{ARGUMENT: \{'The proposed items have a significant total importance to me, and the inclusion of C81 adds diversity to the set. However, I would like to explore the possibility of adding or swapping items to further increase the total importance. I will consider the effort and importance of each item to determine the best course of action.'\}} \\
\\ 
\texttt{PROPOSAL: \{'A08', 'A27', 'B12', 'A16', 'B60', 'C81', 'A84'\}} \\
            }
        }
    }
    & \\ \\

    \theutterance \stepcounter{utterance}  
    & \multicolumn{4}{p{0.6\linewidth}}{
        \cellcolor[rgb]{0.9,0.9,0.9}{
            \makecell[{{p{\linewidth}}}]{
                \texttt{\tiny{[P1$\rangle$GM]}}
                \texttt{STRATEGIC REASONING: \{'The other participant has added A84 to the proposal. I need to evaluate the new proposal and decide whether to accept it or make a counter{-}proposal. The total effort of the new proposal is 607 + 106 + 946 + 645 + 324 + 601 + 136 = 3365, which is within the LIMIT. I will calculate the total importance of the new proposal and compare it to the previous proposals. I will also consider the importance of the added item A84 and its impact on the overall set.'\}} \\
\\ 
\texttt{REFUSE: \{'A08', 'A27', 'B12', 'A16', 'B60', 'C81', 'A84'\}} \\
\\ 
\texttt{ARGUMENT: \{'While the added item A84 does increase the total importance, I believe it may not be the most efficient choice. The effort of A84 is relatively low, but its importance is also moderate. I would like to explore other options that could potentially increase the total importance more significantly. I refuse the current proposal and will make a new proposal that includes items with high importance to me.'\}} \\
\\ 
\texttt{PROPOSAL: \{'A08', 'A27', 'B12', 'A16', 'B60', 'C81', 'A30'\}} \\
            }
        }
    }
    & & \\ \\

    \theutterance \stepcounter{utterance}  
    & & \multicolumn{4}{p{0.6\linewidth}}{
        \cellcolor[rgb]{0.9,0.9,0.9}{
            \makecell[{{p{\linewidth}}}]{
                \texttt{\tiny{[GM$\rangle$P2]}}
                \texttt{REFUSE: \{'A08', 'A27', 'B12', 'A16', 'B60', 'C81', 'A84'\}} \\
\\ 
\texttt{ARGUMENT: \{'While the added item A84 does increase the total importance, I believe it may not be the most efficient choice. The effort of A84 is relatively low, but its importance is also moderate. I would like to explore other options that could potentially increase the total importance more significantly. I refuse the current proposal and will make a new proposal that includes items with high importance to me.'\}} \\
\\ 
\texttt{PROPOSAL: \{'A08', 'A27', 'B12', 'A16', 'B60', 'C81', 'A30'\}} \\
            }
        }
    }
    & \\ \\

    \theutterance \stepcounter{utterance}  
    & & & \multicolumn{4}{p{0.6\linewidth}}{
        \cellcolor[rgb]{0.9,0.9,0.9}{
            \makecell[{{p{\linewidth}}}]{
                \texttt{\tiny{[GM$\langle$P2]}}
                \texttt{STRATEGIC REASONING: \{'The other player has refused my previous proposal and made a new proposal that includes item A30. I will calculate the total effort and importance of this new proposal to determine its value to me. The total effort of the proposed items is 607 + 106 + 946 + 645 + 324 + 601 + 213 = 3442. The total importance of these items to me is 868 + 822 + 915 + 856 + 808 + 713 + 739 = 5621. This proposal has a high total importance to me, and the inclusion of A30 adds significant value. I will consider accepting this proposal, but I will also evaluate the potential benefits of making a counter{-}proposal to further increase the total importance.'\}} \\
\\ 
\texttt{ARGUMENT: \{'The new proposal includes item A30, which has a high importance to me. The total effort of the proposed items is still within the LIMIT, and the total importance is significant. I will consider accepting this proposal, but I would like to confirm that it is the best possible option.'\}} \\
\\ 
\texttt{PROPOSAL: \{'A08', 'A27', 'B12', 'A16', 'B60', 'C81', 'A30'\}} \\
            }
        }
    }
     \\ \\

    \theutterance \stepcounter{utterance}  
    & & \multicolumn{4}{p{0.6\linewidth}}{
        \cellcolor[rgb]{0.9,0.9,0.9}{
            \makecell[{{p{\linewidth}}}]{
                \texttt{\tiny{[P1$\langle$GM]}}
                \texttt{ARGUMENT: \{'The new proposal includes item A30, which has a high importance to me. The total effort of the proposed items is still within the LIMIT, and the total importance is significant. I will consider accepting this proposal, but I would like to confirm that it is the best possible option.'\}} \\
\\ 
\texttt{PROPOSAL: \{'A08', 'A27', 'B12', 'A16', 'B60', 'C81', 'A30'\}} \\
            }
        }
    }
    & \\ \\

    \theutterance \stepcounter{utterance}  
    & \multicolumn{4}{p{0.6\linewidth}}{
        \cellcolor[rgb]{0.9,0.9,0.9}{
            \makecell[{{p{\linewidth}}}]{
                \texttt{\tiny{[P1$\rangle$GM]}}
                \texttt{STRATEGIC REASONING: \{'The other participant has confirmed the proposal I made earlier, which includes A30. I need to evaluate whether this proposal is indeed the best possible option. The total effort of the proposal is 607 + 106 + 946 + 645 + 324 + 601 + 213 = 3542, which is within the LIMIT. I will review the importance of each item in the proposal and consider whether any other items could potentially increase the total importance.'\}} \\
\\ 
\texttt{ARGUMENT: \{'I have reviewed the proposal and believe it is a strong option. The inclusion of A30 adds significant importance, and the total effort is still within the LIMIT. I am willing to accept this proposal as it seems to be a good balance of importance and effort.'\}} \\
\\ 
\texttt{AGREE: \{'A08', 'A27', 'B12', 'A16', 'B60', 'C81', 'A30'\}} \\
            }
        }
    }
    & & \\ \\

    \theutterance \stepcounter{utterance}  
    & & & \multicolumn{2}{p{0.3\linewidth}}{
        \cellcolor[rgb]{0.9,0.9,0.9}{
            \makecell[{{p{\linewidth}}}]{
                \texttt{\tiny{[GM$|$GM]}}
                \texttt{game successful} \\
            }
        }
    }
    & & \\ \\

    \theutterance \stepcounter{utterance}  
    & & & \multicolumn{2}{p{0.3\linewidth}}{
        \cellcolor[rgb]{0.9,0.9,0.9}{
            \makecell[{{p{\linewidth}}}]{
                \texttt{\tiny{[GM$|$GM]}}
                \texttt{end game} \\
            }
        }
    }
    & & \\ \\

\end{supertabular}
}

\end{document}
